\documentclass[10pt,a4paper]{article}
\usepackage[margin=20mm,headheight=15pt,headsep=7mm,footskip=10mm]{geometry}
\usepackage{fontspec,amsmath,amssymb,mathtools,microtype}
\setmainfont{lmroman10-regular.otf}[BoldFont=lmroman10-bold.otf,ItalicFont=lmroman10-italic.otf,BoldItalicFont=lmroman10-bolditalic.otf]
\setsansfont{lmsans10-regular.otf}[BoldFont=lmsans10-bold.otf,ItalicFont=lmsans10-oblique.otf,BoldItalicFont=lmsans10-boldoblique.otf]
\setmonofont[Scale=.86]{lmmono10-regular.otf}[ItalicFont=lmmono10-italic.otf]
\usepackage{ragged2e}\usepackage{xcolor,booktabs,longtable,array,graphicx,enumitem,fancyhdr,needspace,titlesec,hyperref,xurl}
\definecolor{AtlasNavy}{HTML}{18334C}\definecolor{AtlasTeal}{HTML}{116D73}\definecolor{AtlasGray}{HTML}{52616D}
\hypersetup{colorlinks=true,linkcolor=AtlasTeal,urlcolor=AtlasTeal,pdfcreator={Complexity Atlas},unicode=true}
\urlstyle{same}\setlength{\parindent}{0pt}\setlength{\parskip}{4.8pt}\setlength{\emergencystretch}{3em}
\renewcommand{\baselinestretch}{1.08}\setcounter{secnumdepth}{0}\setlist{nosep,leftmargin=1.5em}
\titleformat{\section}{\Large\bfseries\sffamily\color{AtlasNavy}\raggedright}{}{0em}{}
\titleformat{\subsection}{\large\bfseries\sffamily\color{AtlasTeal}\raggedright}{}{0em}{}
\titleformat{\subsubsection}{\normalsize\bfseries\color{AtlasNavy}\raggedright}{}{0em}{}
\titlespacing*{\section}{0pt}{14pt}{6pt}\titlespacing*{\subsection}{0pt}{10pt}{5pt}
\pagestyle{fancy}\fancyhf{}\fancyhead[L]{\small\sffamily\color{AtlasGray}COMPLEXITY ATLAS}
\fancyhead[R]{\small\sffamily\color{AtlasGray}ENGLISH EDITION}\fancyfoot[C]{\small\thepage}\renewcommand{\headrulewidth}{.3pt}
\newcommand{\AtlasMath}[1]{\par\begingroup\centering\setbox0=\hbox{$\displaystyle #1$}\ifdim\wd0>1.4\linewidth\typeout{ATLAS WIDE FORMULA: \the\wd0}\fi\ifdim\wd0>\linewidth\resizebox{\linewidth}{!}{\box0}\else\box0\fi\par\endgroup}
\newcommand{\AtlasInlineMath}[1]{\begingroup\setbox0=\hbox{$#1$}\ifdim\wd0>.4\linewidth\endgroup\AtlasMath{#1}\else\endgroup$#1$\fi}
\pdfstringdefDisableCommands{\def\AtlasInlineMath#1{#1}}
\DeclareMathOperator{\polylog}{polylog}\DeclareMathOperator{\dist}{dist}\DeclareMathOperator{\diam}{diam}
\DeclareMathOperator{\rank}{rank}\DeclareMathOperator{\tr}{tr}\DeclareMathOperator{\diag}{diag}
\DeclareMathOperator{\prox}{prox}\DeclareMathOperator{\KL}{KL}\DeclareMathOperator{\Var}{Var}
\providecommand{\eps}{\varepsilon}\providecommand{\R}{\mathbb R}\providecommand{\E}{\mathbb E}
\providecommand{\mathbbm}[1]{\mathbf{#1}}\providecommand{\mathds}[1]{\mathbf{#1}}
\providecommand{\llbracket}{[\![}\providecommand{\rrbracket}{]\!]}
\begin{document}
\raggedright

{\LARGE\bfseries\sffamily\color{AtlasNavy} Minimization of a finite maximum of convex functions: public ball; selected component values and gradients\par}\vspace{5pt}


\textbf{Setting ID.} \texttt{SO\_\allowbreak{}FO\_\allowbreak{}F04\_\allowbreak{}ball}



Single-objective optimization



Find an approximate global minimizer, with accuracy measured by suboptimality of the complete objective.



\textbf{Status.} Logarithmic gap; Other quantified gap



\textbf{Gap classes present.} Logarithmic gap; Other quantified gap



\textbf{TASK tags.} formulation: Global minimization; task family: Single objective



\textbf{MODEL tags.} convexity: Convex; convexity: Convex components; geometry: Euclidean ball; geometry: Full space; smoothness: Component smoothness; smoothness: Lipschitz function; structure: Finite max; structure: Unknown component offsets



\textbf{ORACLE tags.} access order: First-order interface; feedback quality: Exact; interface family: Component value gradient; operation: Component value gradient; public information: Hidden objects explicitly declared; public information: Known upper bound; public information: Known exact value; public information: Public inputs explicitly declared



\textbf{ALGORITHMS tags.} adaptivity: Adaptive; computation: Internal arithmetic uncharged; randomization: Deterministic; randomization: Private randomness allowed; restriction: No additional span restriction explicitly imposed



\textbf{OUTPUT tags.} certificate: No additional output certificate required; feasibility: Required; object: Point; postprocessing: Unqueried output explicitly allowed



\textbf{METRIC tags.} criterion: Objective gap



\textbf{GUARANTEE tags.} aggregation: Per instance; mode: Deterministic; mode: Expectation; mode: Probability; quantifier: Algorithm before instance; randomness: Algorithm private seed



\textbf{RESOURCE tags.} budget: Fixed worst-case budget; measure: Component calls; precision: Bit precision not counted



\needspace{7\baselineskip}\section{1. Setting specification}


\needspace{4\baselineskip}\subsection{Task}


Find an approximate global minimizer, with accuracy measured by suboptimality of the complete objective.




\AtlasMath{\min_{x\in Q}F(x)}




The complete objective is \AtlasInlineMath{F(x)=\max_{1\le i\le n}f_i(x)}. It may be nonsmooth even when all components are smooth.



\needspace{4\baselineskip}\subsection{Admissible model}


The \AtlasInlineMath{n} unknown components are convex and differentiable in a neighborhood of the public ball \AtlasInlineMath{Q}, with \AtlasInlineMath{\|\nabla f_i\|\le L} and \AtlasInlineMath{\beta}-Lipschitz gradients. Their constant offsets are unknown. The maximum \AtlasInlineMath{F=\max_i f_i} is generally nonsmooth.



Feasible geometry: \AtlasInlineMath{Q=\{x\in\mathbb{R}^d:\|x\|_2\le R\}}.



Unless another geometry is specified, vectors use the Euclidean norm \AtlasInlineMath{\|\cdot\|_2}, and matrices and higher derivatives use the corresponding operator norm. The assumptions hold uniformly throughout the declared domain.



Individual smoothness means \AtlasInlineMath{\|\nabla f_i(u)-\nabla f_i(v)\|_2\le\beta\|u-v\|_2} simultaneously for every \AtlasInlineMath{i\in\{1,\ldots,n\}} and every pair \AtlasInlineMath{u,v} in the declared model domain.



Admissible parameter domain:


\AtlasMath{\begin{gathered}\varepsilon > 0,\quad d \ge 1,\quad n \ge 1,\quad L > 0\\R > 0,\quad \beta \ge 0\end{gathered}}




\needspace{4\baselineskip}\subsection{Oracle interface}


Public numerical inputs: \AtlasInlineMath{d,\ n,\ \varepsilon}.



Public upper bounds: \AtlasInlineMath{L,\ \beta,\ R}.



The unknown objects are the component functions \AtlasInlineMath{f_1,\ldots,f_n}.



For each query, the algorithm chooses \AtlasInlineMath{i\in\{1,\ldots,n\}} and \AtlasInlineMath{x\in Q} and receives the following information about that component.




\AtlasMath{(f_i(x),\nabla f_i(x))}




All responses come from one fixed unknown instance. The algorithm adaptively chooses query points and any indices permitted by the interface.



Exact Euclidean projection onto the public domain; its internal computation is free



\needspace{4\baselineskip}\subsection{Output}


Return a point \AtlasInlineMath{\widehat x\in Q}.



The output may be unqueried or constructed from the query history, provided it is feasible. No additional function value or accuracy certificate is required.



\needspace{4\baselineskip}\subsection{Loss or accuracy criterion}



\AtlasMath{F(\widehat{x})-\min_{u\in Q}F(u)}




\AtlasInlineMath{F^\star=\inf_{x\in Q}F(x)}. The existence, compactness, or initial-suboptimality assumptions ensure that the optimal value in this comparison is finite.



\needspace{4\baselineskip}\subsection{Resource accounting}


One component value-gradient query at an algorithm-selected pair \AtlasInlineMath{(i,x)} costs one unit; querying all \AtlasInlineMath{n} components costs \AtlasInlineMath{n}. Exact internal solution of a public convex model is uncharged.



Cost is bounded uniformly on every execution path. All calls used for initialization, repetition, batching, or candidate verification are counted.



Real arithmetic, precision, and storage applied to already available information are uncharged. Public geometric operations are limited to those explicitly granted by the oracle. Query complexity alone does not imply arithmetic, memory, or communication complexity.



\needspace{4\baselineskip}\subsection{Algorithm views and success guarantees}


\needspace{4\baselineskip}\subsubsection{Deterministic / Pointwise guarantee}


The algorithm uses no private random numbers. Query and stopping decisions may depend adaptively on the observed feedback.



No additional span, zero-respecting, linear-iteration, or memory restriction is imposed.



Interaction is adaptive to the feedback history. Every actual response is charged separately, including responses in parallel or batch queries. Internal real arithmetic, storage, and postprocessing are uncharged.



Select an algorithm using the public parameters, and then require the following for every admissible instance and response rule:




\AtlasMath{\forall I\,\forall\mathcal{O}\in\mathcal{O}(I):\ \ell(I,\widehat{x})\le\varepsilon}




The error guarantee holds for every instance; it is not averaged over an input distribution.



\needspace{4\baselineskip}\subsubsection{Randomized / Expected loss}


Private randomization is allowed. Queries, stopping decisions, and postprocessing may adapt to observed feedback.



No additional span, zero-respecting, linear-iteration, or memory restriction is imposed.



Interaction is adaptive to the feedback history. Every actual response is charged separately, including responses in parallel or batch queries. Internal real arithmetic, storage, and postprocessing are uncharged.



Select an algorithm using the public parameters, and then require the following for every admissible instance and response rule:




\AtlasMath{\forall I\,\forall\mathcal{O}\in\mathcal{O}(I):\ \mathbb{E}_{A,\mathcal{O}}\ell(I,\widehat{x})\le\varepsilon}




The probability space includes the algorithm's private randomness. The instance and its admissible oracle law are fixed before randomness is averaged.



\needspace{4\baselineskip}\subsubsection{Randomized / 2/3 success / normalized global regularity}


Queries, stopping decisions, and the output are measurable functions of the preceding transcript and an instance-independent private random seed. The unknown instance is fixed before the seed is realized. No span, support, zero-respecting, or memory restriction is imposed.



Internal exact-real computation and storage using public inputs and received information are uncharged. The stated budget must hold on every execution path.




\AtlasMath{\forall(f_i):\ \Pr(F(\widehat x)-\min_{x\in B_2^d}F(x)\le\varepsilon)\ge2/3}




The guarantee holds for every fixed admissible component tuple. Probability is over private algorithmic randomness, and every execution uses at most the prescribed number of component queries.



Model: There are \AtlasInlineMath{n} hidden globally convex differentiable functions \AtlasInlineMath{f_i:\mathbb R^d\to\mathbb R}. For all \AtlasInlineMath{x,y\in\mathbb R^d} and all \AtlasInlineMath{i}, \AtlasInlineMath{|f_i(x)-f_i(y)|\le\|x-y\|_2} and \AtlasInlineMath{\|\nabla f_i(x)-\nabla f_i(y)\|_2\le\|x-y\|_2}. Thus \AtlasInlineMath{L=\beta=R=1}. Component offsets are unrestricted and unknown.



Model: The public feasible set is \AtlasInlineMath{Q=B_2^d}. Query points and the final output must belong to \AtlasInlineMath{Q}, although the regularity assumptions are global. The objective is \AtlasInlineMath{F(x)=\max_{1\le i\le n}f_i(x)}.



\textbf{Budget convention.} fixed\_worst\_case



\needspace{7\baselineskip}\section{2. Notation and parameters}


\AtlasInlineMath{x,\widehat{x}}: Query points and final output; their permitted spaces are specified separately by the oracle and output contracts \textbf{Domain:} \AtlasInlineMath{\widehat{x}\in Q}.



\AtlasInlineMath{f_i,F}: Unknown components and finite-maximum objective \textbf{Domain:} \AtlasInlineMath{i\in\{1,\ldots,n\}}.




\AtlasMath{F(x)=\max_{1\le i\le n}f_i(x)}




\AtlasInlineMath{\nabla f_i}: Full gradient of differentiable component \AtlasInlineMath{i}



\AtlasInlineMath{F^\star}: Optimal value of the complete task; the oracle need not return it




\AtlasMath{F^\star=\inf_{x\in Q}F(x)}




\AtlasInlineMath{Q}: Feasible set of the task




\AtlasMath{Q=\{x\in\mathbb{R}^d:\|x\|_2\le R\}}




\AtlasInlineMath{\mathbb{R}^d}: Real vector space with \AtlasInlineMath{d} coordinates



\AtlasInlineMath{\|\cdot\|_2,\langle\cdot,\cdot\rangle}: Euclidean norm and inner product



\AtlasInlineMath{d}: Ambient dimension; a positive integer tracked explicitly. \textbf{Domain:} \AtlasInlineMath{d\in\mathbb{N}_{\ge1}}.



\AtlasInlineMath{n}: Public number of components in the finite sum or finite maximum. \textbf{Domain:} \AtlasInlineMath{n\in\mathbb{N}_{\ge1}}.



\AtlasInlineMath{L}: Public upper bound on the Lipschitz constant of the objective or components. \textbf{Domain:} \AtlasInlineMath{L>0}.



\AtlasInlineMath{\beta}: Public upper bound on the gradient smoothness constant. Any mean-square interpretation is specified in the model. \textbf{Domain:} \AtlasInlineMath{\beta\ge0}.



\AtlasInlineMath{R}: Radius of the public feasible ball \textbf{Domain:} \AtlasInlineMath{R>0}.



\AtlasInlineMath{\varepsilon}: Target tolerance for the underlying error defined by the accuracy measure \textbf{Domain:} \AtlasInlineMath{\varepsilon>0}.



\AtlasInlineMath{c}: Sufficiently small universal positive constant, potentially different between theorems; it is not a numerical threshold. \textbf{Domain:} \AtlasInlineMath{0<c<1}.



\AtlasInlineMath{C}: Sufficiently large universal constant, independent of the tracked parameters; no numerical value is specified. \textbf{Domain:} \AtlasInlineMath{C>1}.



\AtlasInlineMath{A_\varepsilon}: Dimensionless quantity \AtlasInlineMath{LR/\epsilon}, used only to shorten formulas.




\AtlasMath{A_\varepsilon=\frac{L\,R}{\varepsilon}}




\AtlasInlineMath{B_\varepsilon}: Dimensionless quantity \AtlasInlineMath{\beta R^2/\epsilon}; not an independent parameter.




\AtlasMath{B_\varepsilon=\frac{\beta\,{R}^{2}}{\varepsilon}}




\AtlasInlineMath{\ell}: The underlying loss or error defined by the accuracy measure; its expected, probabilistic, or worst-case aggregation is specified by the success contract



\AtlasInlineMath{I,\mathcal{O},A}: Admissible instance, a response rule consistent with that instance, and an algorithm; the algorithm is selected before the unknown instance



\AtlasInlineMath{O,\Omega,\Theta,\widetilde{O},\widetilde\Omega}: Notation for asymptotic upper and lower bounds, constant-factor matching rates, and explicitly permitted logarithmic factors. Hidden-constant dependence is specified in each regime.



\AtlasInlineMath{\log}: Natural logarithm; changing to another fixed base affects only a constant



\AtlasInlineMath{\mathbb{E}}: Expectation over the declared randomness after fixing the instance and admissible response law



\needspace{12\baselineskip}\section{3. Complexity bracket}


\needspace{4\baselineskip}\subsection{Compared quantities}


\needspace{4\baselineskip}\subsubsection{Q1 - target accuracy cost}


\textbf{View.} Deterministic / Pointwise guarantee



\AtlasMath{C^\star_{v,\pi}(\varepsilon;\theta)=\inf\{B\ge0:\exists\mathcal A\in\mathfrak A_v,\ C_\pi(\mathcal A,I)\le B\ \text{and}\ G_v(\mathcal A,I;\varepsilon)\ \forall I\}}


\textbf{Definition.} One component value-gradient query at an algorithm-selected pair \AtlasInlineMath{(i,x)} costs one unit; querying all \AtlasInlineMath{n} components costs \AtlasInlineMath{n}. Exact internal solution of a public convex model is uncharged.



\textbf{Scope.} declared\_joint\_parameter\_window



\textbf{Optimization direction.} minimize



Here v is the selected algorithm view, G\_v its complete success guarantee, I an admissible fixed instance, theta the remaining parameters, and C\_pi the declared resource projection. C\_pi uses the view's fixed-path or expected-work budget convention, including initialization. No change of guarantee or scalarization is implicit.



\needspace{4\baselineskip}\subsubsection{Q2 - target accuracy cost}


\textbf{View.} Deterministic / Pointwise guarantee



\AtlasMath{C^\star_{v,\pi}(\varepsilon;\theta)=\inf\{B\ge0:\exists\mathcal A\in\mathfrak A_v,\ C_\pi(\mathcal A,I)\le B\ \text{and}\ G_v(\mathcal A,I;\varepsilon)\ \forall I\}}


\textbf{Definition.} One component value-gradient query at an algorithm-selected pair \AtlasInlineMath{(i,x)} costs one unit; querying all \AtlasInlineMath{n} components costs \AtlasInlineMath{n}. Exact internal solution of a public convex model is uncharged.



\textbf{Scope.} declared\_joint\_parameter\_window



\textbf{Optimization direction.} minimize



Here v is the selected algorithm view, G\_v its complete success guarantee, I an admissible fixed instance, theta the remaining parameters, and C\_pi the declared resource projection. C\_pi uses the view's fixed-path or expected-work budget convention, including initialization. No change of guarantee or scalarization is implicit.



\needspace{4\baselineskip}\subsubsection{Q3 - target accuracy cost}


\textbf{View.} Randomized / Expected loss



\AtlasMath{C^\star_{v,\pi}(\varepsilon;\theta)=\inf\{B\ge0:\exists\mathcal A\in\mathfrak A_v,\ C_\pi(\mathcal A,I)\le B\ \text{and}\ G_v(\mathcal A,I;\varepsilon)\ \forall I\}}


\textbf{Definition.} One component value-gradient query at an algorithm-selected pair \AtlasInlineMath{(i,x)} costs one unit; querying all \AtlasInlineMath{n} components costs \AtlasInlineMath{n}. Exact internal solution of a public convex model is uncharged.



\textbf{Scope.} declared\_joint\_parameter\_window



\textbf{Optimization direction.} minimize



Here v is the selected algorithm view, G\_v its complete success guarantee, I an admissible fixed instance, theta the remaining parameters, and C\_pi the declared resource projection. C\_pi uses the view's fixed-path or expected-work budget convention, including initialization. No change of guarantee or scalarization is implicit.



\needspace{4\baselineskip}\subsubsection{Q4 - target accuracy cost}


\textbf{View.} Randomized / Expected loss



\AtlasMath{C^\star_{v,\pi}(\varepsilon;\theta)=\inf\{B\ge0:\exists\mathcal A\in\mathfrak A_v,\ C_\pi(\mathcal A,I)\le B\ \text{and}\ G_v(\mathcal A,I;\varepsilon)\ \forall I\}}


\textbf{Definition.} One component value-gradient query at an algorithm-selected pair \AtlasInlineMath{(i,x)} costs one unit; querying all \AtlasInlineMath{n} components costs \AtlasInlineMath{n}. Exact internal solution of a public convex model is uncharged.



\textbf{Scope.} declared\_joint\_parameter\_window



\textbf{Optimization direction.} minimize



Here v is the selected algorithm view, G\_v its complete success guarantee, I an admissible fixed instance, theta the remaining parameters, and C\_pi the declared resource projection. C\_pi uses the view's fixed-path or expected-work budget convention, including initialization. No change of guarantee or scalarization is implicit.



\needspace{4\baselineskip}\subsubsection{Q5 - minimax metric resource quantity}


\textbf{View.} Randomized / 2/3 success / normalized global regularity



\AtlasMath{Q(n,\varepsilon)=\inf\{T:\forall d\ge1\ \exists A_d\text{ with cap }T\text{ and success }\ge2/3\}}


\textbf{Scope.} The union of the explicit parameter and dimension domains of the linked comparison regimes; each bound retains its own applicability conditions.



\textbf{Optimization direction.} minimize



\textbf{Resource contract.} expression: \AtlasInlineMath{N_{\mathrm{component\ value+gradient}}}; description: All charged oracle operations are counted, including initialization and verification. The budget is a worst-case cap unless the resource expression explicitly specifies another cost.



\needspace{4\baselineskip}\subsection{R1 - Sufficient dimension for the lower bound against arbitrary randomized algorithms}


\textbf{Region.} \AtlasInlineMath{\varepsilon\le c\min\{LR,\beta R^2\},\quad d\ge C\min\{A_\varepsilon^2,B_\varepsilon\}\log(2+n\min\{A_\varepsilon,B_\varepsilon\})}



\textbf{Views.} Deterministic / Pointwise guarantee



\textbf{Quantity.} Q1



\textbf{Resource.} One component value-gradient query at an algorithm-selected pair \AtlasInlineMath{(i,x)} costs one unit; querying all \AtlasInlineMath{n} components costs \AtlasInlineMath{n}. Exact internal solution of a public convex model is uncharged.



\textbf{Status.} Other quantified gap



\textbf{Lower bound.} B4, B5, B6



\AtlasMath{c[n\min\{A_\varepsilon^{2/3},B_\varepsilon^{1/3}\}+\min\{A_\varepsilon^2,\sqrt{nB_\varepsilon}\}]}


\textbf{Upper bound.} B1, B2, B3



\AtlasMath{\min\{Cn(1+\sqrt{B_\varepsilon}),Cn(1+A_\varepsilon^2),Cn[1+d\log(2+A_\varepsilon)]\}}


\textbf{Gap or comparison.}
\AtlasMath{\frac{\min\{Cn(1+\sqrt{B_\varepsilon}),Cn(1+A_\varepsilon^2),Cn[1+d\log(2+A_\varepsilon)]\}}{c[n\min\{A_\varepsilon^{2/3},B_\varepsilon^{1/3}\}+\min\{A_\varepsilon^2,\sqrt{nB_\varepsilon}\}]}}




\textbf{Comparison argument.} Compare on the intersection of F04.ball-D.parameter-domain and F04.hd. Retain every polynomial term and the envelope of applicable algorithms. The lower bound requires sufficiently high dimension.



\needspace{4\baselineskip}\subsection{R2 - One-dimensional smooth accuracy lower bound}


\textbf{Region.} \AtlasInlineMath{0<\varepsilon\le c\min\{LR,\beta R^2\}}



\textbf{Views.} Deterministic / Pointwise guarantee



\textbf{Quantity.} Q1



\textbf{Resource.} One component value-gradient query at an algorithm-selected pair \AtlasInlineMath{(i,x)} costs one unit; querying all \AtlasInlineMath{n} components costs \AtlasInlineMath{n}. Exact internal solution of a public convex model is uncharged.



\textbf{Status.} Other quantified gap



\textbf{Lower bound.} B5, B6



\AtlasMath{c\max\{n,\log(2+\min\{LR,\beta R^2\}/\varepsilon)\}}


\textbf{Upper bound.} B1, B2, B3



\AtlasMath{Cn[1+\min\{\sqrt{B_\varepsilon},A_\varepsilon^2,d\log(2+A_\varepsilon)\}]}


\textbf{Gap or comparison.} \AtlasInlineMath{\frac{Cn[1+\min\{\sqrt{B_\varepsilon},A_\varepsilon^2,d\log(2+A_\varepsilon)\}]}{c\max\{n,\log(2+\min\{LR,\beta R^2\}/\varepsilon)\}}}



\textbf{Comparison argument.} Compare on the intersection of F04.ball-D.parameter-domain and F04.log-lower. A logarithmic-accuracy upper bound is available for fixed \AtlasInlineMath{d}, but the joint dependence on \AtlasInlineMath{n,d,} and accuracy is not matched.



\needspace{4\baselineskip}\subsection{R3 - Normalized high-dimensional slice}


\textbf{Region.} \AtlasInlineMath{L=\beta=R=1,\quad\theta\in\Theta_{\rm hd}}



\textbf{Views.} Deterministic / Pointwise guarantee



\textbf{Quantity.} Q2



\textbf{Resource.} One component value-gradient query at an algorithm-selected pair \AtlasInlineMath{(i,x)} costs one unit; querying all \AtlasInlineMath{n} components costs \AtlasInlineMath{n}. Exact internal solution of a public convex model is uncharged.



\textbf{Status.} Optimal



\textbf{Lower bound.} B4



\AtlasMath{c_n\varepsilon^{-1/2}}


\textbf{Upper bound.} B1, B2, B3



\AtlasMath{C_n\varepsilon^{-1/2}}


\textbf{Gap or comparison.} \AtlasInlineMath{O_n(1)}



\textbf{Comparison argument.} Compare on the intersection of F04.ball-D.parameter-domain and F04.normalized. The number of components \AtlasInlineMath{n} is fixed in advance. Constant-factor optimality jointly in \AtlasInlineMath{n} and \AtlasInlineMath{\epsilon} is not asserted.



\needspace{4\baselineskip}\subsection{R4 - Affine-component endpoint}


\textbf{Region.} \AtlasInlineMath{\beta=0,\quad0<\varepsilon\le cLR}



\textbf{Views.} Deterministic / Pointwise guarantee



\textbf{Quantity.} Q1



\textbf{Resource.} One component value-gradient query at an algorithm-selected pair \AtlasInlineMath{(i,x)} costs one unit; querying all \AtlasInlineMath{n} components costs \AtlasInlineMath{n}. Exact internal solution of a public convex model is uncharged.



\textbf{Status.} Optimal



\textbf{Lower bound.} B5



\AtlasMath{cn}


\textbf{Upper bound.} B7



\AtlasMath{n}


\textbf{Gap or comparison.} \AtlasInlineMath{O(1)}



\needspace{4\baselineskip}\subsection{R5 - Return the ball center without queries}


\textbf{Region.} \AtlasInlineMath{\varepsilon\ge LR}



\textbf{Views.} Deterministic / Pointwise guarantee



\textbf{Quantity.} Q1



\textbf{Resource.} One component value-gradient query at an algorithm-selected pair \AtlasInlineMath{(i,x)} costs one unit; querying all \AtlasInlineMath{n} components costs \AtlasInlineMath{n}. Exact internal solution of a public convex model is uncharged.



\textbf{Status.} Optimal



\textbf{Lower bound.} Nonnegativity of the compared quantity



\AtlasMath{0}


\textbf{Upper bound.} B8



\AtlasMath{0}


\needspace{4\baselineskip}\subsection{R6 - Sufficient dimension for the lower bound against arbitrary randomized algorithms}


\textbf{Region.} \AtlasInlineMath{\varepsilon\le c\min\{LR,\beta R^2\},\quad d\ge C\min\{A_\varepsilon^2,B_\varepsilon\}\log(2+n\min\{A_\varepsilon,B_\varepsilon\})}



\textbf{Views.} Randomized / Expected loss



\textbf{Quantity.} Q3



\textbf{Resource.} One component value-gradient query at an algorithm-selected pair \AtlasInlineMath{(i,x)} costs one unit; querying all \AtlasInlineMath{n} components costs \AtlasInlineMath{n}. Exact internal solution of a public convex model is uncharged.



\textbf{Status.} Other quantified gap



\textbf{Lower bound.} B13, B14, B15



\AtlasMath{c[n\min\{A_\varepsilon^{2/3},B_\varepsilon^{1/3}\}+\min\{A_\varepsilon^2,\sqrt{nB_\varepsilon}\}]}


\textbf{Upper bound.} B9, B10, B11, B12



\AtlasMath{\min\{Cn(1+\sqrt{B_\varepsilon}),Cn(1+A_\varepsilon^2),Cn[1+d\log(2+A_\varepsilon)],\widetilde{O}(nB_\varepsilon^{1/3}+A_\varepsilon^2)\}}


\textbf{Gap or comparison.}
\AtlasMath{\frac{\min\{Cn(1+\sqrt{B_\varepsilon}),Cn(1+A_\varepsilon^2),Cn[1+d\log(2+A_\varepsilon)],\widetilde{O}(nB_\varepsilon^{1/3}+A_\varepsilon^2)\}}{c[n\min\{A_\varepsilon^{2/3},B_\varepsilon^{1/3}\}+\min\{A_\varepsilon^2,\sqrt{nB_\varepsilon}\}]}}




\textbf{Comparison argument.} Compare on the intersection of F04.ball-R.parameter-domain and F04.hd. Retain every polynomial term and the envelope of applicable algorithms. The lower bound requires sufficiently high dimension.



\needspace{4\baselineskip}\subsection{R7 - One-dimensional smooth accuracy lower bound}


\textbf{Region.} \AtlasInlineMath{0<\varepsilon\le c\min\{LR,\beta R^2\}}



\textbf{Views.} Randomized / Expected loss



\textbf{Quantity.} Q3



\textbf{Resource.} One component value-gradient query at an algorithm-selected pair \AtlasInlineMath{(i,x)} costs one unit; querying all \AtlasInlineMath{n} components costs \AtlasInlineMath{n}. Exact internal solution of a public convex model is uncharged.



\textbf{Status.} Other quantified gap



\textbf{Lower bound.} B14, B15



\AtlasMath{c\max\{n,\log(2+\min\{LR,\beta R^2\}/\varepsilon)\}}


\textbf{Upper bound.} B9, B10, B11



\AtlasMath{Cn[1+\min\{\sqrt{B_\varepsilon},A_\varepsilon^2,d\log(2+A_\varepsilon)\}]}


\textbf{Gap or comparison.} \AtlasInlineMath{\frac{Cn[1+\min\{\sqrt{B_\varepsilon},A_\varepsilon^2,d\log(2+A_\varepsilon)\}]}{c\max\{n,\log(2+\min\{LR,\beta R^2\}/\varepsilon)\}}}



\textbf{Comparison argument.} Compare on the intersection of F04.ball-R.parameter-domain and F04.log-lower. A logarithmic-accuracy upper bound is available for fixed \AtlasInlineMath{d}, but the joint dependence on \AtlasInlineMath{n,d,} and accuracy is not matched.



\needspace{4\baselineskip}\subsection{R8 - Normalized regime with many components}


\textbf{Region.} \AtlasInlineMath{n\ge C\varepsilon^{-5/3},\quad L=\beta=R=1,\quad\theta\in\Theta_{\rm hd}}



\textbf{Views.} Randomized / Expected loss



\textbf{Quantity.} Q4



\textbf{Resource.} One component value-gradient query at an algorithm-selected pair \AtlasInlineMath{(i,x)} costs one unit; querying all \AtlasInlineMath{n} components costs \AtlasInlineMath{n}. Exact internal solution of a public convex model is uncharged.



\textbf{Status.} Logarithmic gap



\textbf{Lower bound.} B13



\AtlasMath{\Omega(n\varepsilon^{-1/3})}


\textbf{Upper bound.} B9, B12, B11



\AtlasMath{\widetilde{O}(n\varepsilon^{-1/3})}


\textbf{Gap or comparison.} \AtlasInlineMath{\operatorname{polylog}(n,1/\varepsilon)}



\textbf{Comparison argument.} Compare on the intersection of F04.ball-R.parameter-domain and F04.large-n. The term \AtlasInlineMath{n\epsilon^{-1/3}} dominates \AtlasInlineMath{\epsilon^{-2}}. The randomized upper bound does not remove the distinction between deterministic and randomized views.



\needspace{4\baselineskip}\subsection{R9 - Affine-component endpoint}


\textbf{Region.} \AtlasInlineMath{\beta=0,\quad0<\varepsilon\le cLR}



\textbf{Views.} Randomized / Expected loss



\textbf{Quantity.} Q3



\textbf{Resource.} One component value-gradient query at an algorithm-selected pair \AtlasInlineMath{(i,x)} costs one unit; querying all \AtlasInlineMath{n} components costs \AtlasInlineMath{n}. Exact internal solution of a public convex model is uncharged.



\textbf{Status.} Optimal



\textbf{Lower bound.} B14



\AtlasMath{cn}


\textbf{Upper bound.} B16



\AtlasMath{n}


\textbf{Gap or comparison.} \AtlasInlineMath{O(1)}



\needspace{4\baselineskip}\subsection{R10 - Return the ball center without queries}


\textbf{Region.} \AtlasInlineMath{\varepsilon\ge LR}



\textbf{Views.} Randomized / Expected loss



\textbf{Quantity.} Q3



\textbf{Resource.} One component value-gradient query at an algorithm-selected pair \AtlasInlineMath{(i,x)} costs one unit; querying all \AtlasInlineMath{n} components costs \AtlasInlineMath{n}. Exact internal solution of a public convex model is uncharged.



\textbf{Status.} Optimal



\textbf{Lower bound.} Nonnegativity of the compared quantity



\AtlasMath{0}


\textbf{Upper bound.} B17



\AtlasMath{0}


\needspace{4\baselineskip}\subsection{R11 - Normalized global regularity: dimension-unrestricted probability complexity}


\textbf{Region.} \AtlasInlineMath{n\ge1} and \AtlasInlineMath{0<\varepsilon<1}; the complexity takes the worst case over all finite dimensions. The upper bound holds for every finite d. The lower bound is existential in dimension, which may depend on n and epsilon; A2 Proposition prop:lower and Section sec:lower do not give an explicit dimensional formula. The normalized global-regularity view is distinct from assuming regularity only near the ball. The comparison concerns oracle calls, not arithmetic runtime.



\textbf{Views.} Randomized / 2/3 success / normalized global regularity



\textbf{Quantity.} Q5



\textbf{Resource.} All charged oracle operations are counted, including initialization and verification. The budget is a worst-case cap unless the resource expression explicitly specifies another cost.



\textbf{Status.} Logarithmic gap



\textbf{Lower bound.} B19



\AtlasMath{c(n\varepsilon^{-1/3}+\sqrt{n/\varepsilon})}


\textbf{Upper bound.} B18



\AtlasMath{C(n\varepsilon^{-1/3}+\sqrt{n/\varepsilon})\log^6(en/\varepsilon)}


\textbf{Gap or comparison.}



\AtlasMath{O(\log^6(en/\varepsilon))}


\needspace{7\baselineskip}\section{4. Bounds and supporting evidence}


\needspace{4\baselineskip}\subsection{B1 - Accelerated iterative linearization}


\textbf{Admission.} Included in the bracket



\textbf{Direction.} upper



\textbf{View.} Deterministic / Pointwise guarantee



\textbf{Witness or lower-bound scope.} Deterministic witness



\AtlasMath{Cn(1+\sqrt{B_\varepsilon})}


\textbf{Sources.} \href{https://arxiv.org/pdf/2105.01778}{[1]}



\textbf{Conditions.}
\AtlasMath{0<\varepsilon\le cLR}




\textbf{Dimension conditions.} \AtlasInlineMath{0<\varepsilon\le cLR}



\textbf{Oracle correspondence.} Applicable interfaces: selected component values and gradients. Algorithms may use only these interfaces and the auxiliary operations explicitly declared public.



\textbf{Guarantee conversion.} A deterministic error bound for every instance also gives the same expected-error bound. Probability guarantees are used only when an explicit conversion to the target contract is recorded.



\textbf{Output rule.} The output must lie in the declared output space and satisfy its feasibility, support, or atomic-hull restrictions. Permission to return an unqueried point does not remove these restrictions.



\textbf{Resource conversion.} Every primitive response is charged. Sample, component, proximal, and communication costs use the units declared for the target setting. Resource coordinates achieved by one algorithm remain a joint vector; coordinatewise minima from different algorithms are not combined.



\textbf{Source location.} \href{https://arxiv.org/pdf/2105.01778}{[1]} theorem: Theorem3 Eq.(13), PDF p.12; AppendixA.3 andD.3



\needspace{4\baselineskip}\subsection{B2 - Projected subgradient descent after full value-subgradient simulation using \AtlasInlineMath{n} component calls}


\textbf{Admission.} Included in the bracket



\textbf{Direction.} upper



\textbf{View.} Deterministic / Pointwise guarantee



\textbf{Witness or lower-bound scope.} Deterministic witness



\AtlasMath{Cn(1+A_\varepsilon^2)}


\textbf{Sources.} \href{https://arxiv.org/pdf/2311.10886}{[3]} \href{https://arxiv.org/pdf/1405.4980}{[4]} \href{https://www2.isye.gatech.edu/~nemirovs/MOR_AccuracyCertificates.pdf}{[5]} \href{https://arxiv.org/pdf/2609.20701}{[6]}



\textbf{Conditions.}
\AtlasMath{0<\varepsilon\le cLR}




\textbf{Dimension conditions.} \AtlasInlineMath{0<\varepsilon\le cLR}



\textbf{Oracle correspondence.} Applicable interfaces: selected component values and gradients. Algorithms may use only these interfaces and the auxiliary operations explicitly declared public.



\textbf{Guarantee conversion.} A deterministic error bound for every instance also gives the same expected-error bound. Probability guarantees are used only when an explicit conversion to the target contract is recorded.



\textbf{Output rule.} The output must lie in the declared output space and satisfy its feasibility, support, or atomic-hull restrictions. Permission to return an unqueried point does not remove these restrictions.



\textbf{Resource conversion.} Every primitive response is charged. Sample, component, proximal, and communication costs use the units declared for the target setting. Resource coordinates achieved by one algorithm remain a joint vector; coordinatewise minima from different algorithms are not combined.



\textbf{Local argument.} FO\_FINITE\_MAXIMUM, FO\_PROJECTED\_SG, FO\_CERTIFICATE, FO\_NESTED\_FG



\needspace{4\baselineskip}\subsection{B3 - Cutting planes after full value-subgradient simulation using \AtlasInlineMath{n} component calls}


\textbf{Admission.} Included in the bracket



\textbf{Direction.} upper



\textbf{View.} Deterministic / Pointwise guarantee



\textbf{Witness or lower-bound scope.} Deterministic witness



\AtlasMath{Cn[1+d\log(2+A_\varepsilon)]}


\textbf{Sources.} \href{https://www2.isye.gatech.edu/~nemirovs/MOR_AccuracyCertificates.pdf}{[2]}



\textbf{Conditions.}
\AtlasMath{0<\varepsilon\le cLR}




\textbf{Dimension conditions.} \AtlasInlineMath{0<\varepsilon\le cLR}



\textbf{Oracle correspondence.} Applicable interfaces: selected component values and gradients. Algorithms may use only these interfaces and the auxiliary operations explicitly declared public.



\textbf{Guarantee conversion.} A deterministic error bound for every instance also gives the same expected-error bound. Probability guarantees are used only when an explicit conversion to the target contract is recorded.



\textbf{Output rule.} The output must lie in the declared output space and satisfy its feasibility, support, or atomic-hull restrictions. Permission to return an unqueried point does not remove these restrictions.



\textbf{Resource conversion.} Every primitive response is charged. Sample, component, proximal, and communication costs use the units declared for the target setting. Resource coordinates achieved by one algorithm remain a joint vector; coordinatewise minima from different algorithms are not combined.



\textbf{Source location.} \href{https://www2.isye.gatech.edu/~nemirovs/MOR_AccuracyCertificates.pdf}{[2]} theorem: Accuracy certificates; combined with F04 FSG simulation



\needspace{4\baselineskip}\subsection{B4 - Hard instances for arbitrary randomized component value-gradient algorithms}


\textbf{Admission.} Included in the bracket



\textbf{Direction.} lower



\textbf{View.} Deterministic / Pointwise guarantee



\textbf{Witness or lower-bound scope.} All randomized adaptive algorithms



\AtlasMath{c[n\min\{A_\varepsilon^{2/3},B_\varepsilon^{1/3}\}+\min\{A_\varepsilon^2,\sqrt{nB_\varepsilon}\}]}


\textbf{Sources.} \href{https://arxiv.org/pdf/2105.01778}{[1]}



\textbf{Conditions.}
\AtlasMath{\varepsilon\le c\min\{LR,\beta R^2\},\quad d\ge C\min\{A_\varepsilon^2,B_\varepsilon\}\log(2+n\min\{A_\varepsilon,B_\varepsilon\})}




\textbf{Dimension conditions.} \AtlasInlineMath{\varepsilon\le c\min\{LR,\beta R^2\},\quad d\ge C\min\{A_\varepsilon^2,B_\varepsilon\}\log(2+n\min\{A_\varepsilon,B_\varepsilon\})}



\textbf{Oracle correspondence.} Applicable interfaces: selected component values and gradients. Lower bounds transfer only through the explicitly recorded simulations from a stronger interface to a weaker one.



\textbf{Guarantee conversion.} The hard-instance quantifiers cover the algorithm class explicitly specified for this bound. Each target setting retains its expected-error or success-probability criterion; a deterministic lower bound is not promoted to a randomized lower bound. Converting the source hard event to an expected-error statement retains the constant tolerance rescaling. Expected running costs additionally require a stopping-time truncation with its constant loss.



\textbf{Output rule.} The output must lie in the declared output space and satisfy its feasibility, support, or atomic-hull restrictions. Permission to return an unqueried point does not remove these restrictions.



\textbf{Resource conversion.} Every primitive response is charged. Sample, component, proximal, and communication costs use the units declared for the target setting. Resource coordinates achieved by one algorithm remain a joint vector; coordinatewise minima from different algorithms are not combined.



\textbf{Source location.} \href{https://arxiv.org/pdf/2105.01778}{[1]} theorem: Theorem3 Eq.(13), PDF p.12; AppendixA.3 andD.3



\needspace{4\baselineskip}\subsection{B5 - Hidden index in arbitrary dimension}


\textbf{Admission.} Included in the bracket



\textbf{Direction.} lower



\textbf{View.} Deterministic / Pointwise guarantee



\textbf{Witness or lower-bound scope.} All randomized adaptive algorithms



\AtlasMath{cn}


\textbf{Sources.} \href{https://arxiv.org/pdf/2311.10886}{[3]} \href{https://arxiv.org/pdf/1405.4980}{[4]} \href{https://www2.isye.gatech.edu/~nemirovs/MOR_AccuracyCertificates.pdf}{[5]} \href{https://arxiv.org/pdf/2609.20701}{[6]}



\textbf{Conditions.}
\AtlasMath{0<\varepsilon\le cLR}




\textbf{Dimension conditions.} \AtlasInlineMath{0<\varepsilon\le cLR}



\textbf{Oracle correspondence.} Applicable interfaces: selected component values and gradients. Lower bounds transfer only through the explicitly recorded simulations from a stronger interface to a weaker one.



\textbf{Guarantee conversion.} The hard-instance quantifiers cover the algorithm class explicitly specified for this bound. Each target setting retains its expected-error or success-probability criterion; a deterministic lower bound is not promoted to a randomized lower bound.



\textbf{Output rule.} The output must lie in the declared output space and satisfy its feasibility, support, or atomic-hull restrictions. Permission to return an unqueried point does not remove these restrictions.



\textbf{Resource conversion.} Every primitive response is charged. Sample, component, proximal, and communication costs use the units declared for the target setting. Resource coordinates achieved by one algorithm remain a joint vector; coordinatewise minima from different algorithms are not combined.



\textbf{Local argument.} FO\_FINITE\_MAXIMUM, FO\_PROJECTED\_SG, FO\_CERTIFICATE, FO\_NESTED\_FG



\needspace{4\baselineskip}\subsection{B6 - Replication of a one-dimensional smooth hard objective}


\textbf{Admission.} Included in the bracket



\textbf{Direction.} lower



\textbf{View.} Deterministic / Pointwise guarantee



\textbf{Witness or lower-bound scope.} All randomized adaptive algorithms



\AtlasMath{c\log(2+\min\{LR,\beta R^2\}/\varepsilon)}


\textbf{Sources.} \href{https://arxiv.org/pdf/2311.10886}{[3]} \href{https://arxiv.org/pdf/1405.4980}{[4]} \href{https://www2.isye.gatech.edu/~nemirovs/MOR_AccuracyCertificates.pdf}{[5]} \href{https://arxiv.org/pdf/2609.20701}{[6]}



\textbf{Conditions.}
\AtlasMath{0<\varepsilon\le c\min\{LR,\beta R^2\}}




\textbf{Dimension conditions.} \AtlasInlineMath{0<\varepsilon\le c\min\{LR,\beta R^2\}}



\textbf{Oracle correspondence.} Applicable interfaces: selected component values and gradients. Lower bounds transfer only through the explicitly recorded simulations from a stronger interface to a weaker one.



\textbf{Guarantee conversion.} The hard-instance quantifiers cover the algorithm class explicitly specified for this bound. Each target setting retains its expected-error or success-probability criterion; a deterministic lower bound is not promoted to a randomized lower bound.



\textbf{Output rule.} The output must lie in the declared output space and satisfy its feasibility, support, or atomic-hull restrictions. Permission to return an unqueried point does not remove these restrictions.



\textbf{Resource conversion.} Every primitive response is charged. Sample, component, proximal, and communication costs use the units declared for the target setting. Resource coordinates achieved by one algorithm remain a joint vector; coordinatewise minima from different algorithms are not combined.



\textbf{Local argument.} FO\_FINITE\_MAXIMUM, FO\_PROJECTED\_SG, FO\_CERTIFICATE, FO\_NESTED\_FG



\needspace{4\baselineskip}\subsection{B7 - Recover all affine component formulas}


\textbf{Admission.} Included in the bracket



\textbf{Direction.} upper



\textbf{View.} Deterministic / Pointwise guarantee



\textbf{Witness or lower-bound scope.} Deterministic witness



\AtlasMath{n}


\textbf{Sources.} \href{https://arxiv.org/pdf/2311.10886}{[3]} \href{https://arxiv.org/pdf/1405.4980}{[4]} \href{https://www2.isye.gatech.edu/~nemirovs/MOR_AccuracyCertificates.pdf}{[5]} \href{https://arxiv.org/pdf/2609.20701}{[6]}



\textbf{Conditions.}
\AtlasMath{\beta=0,\quad0<\varepsilon\le cLR}




\textbf{Dimension conditions.} \AtlasInlineMath{\beta=0,\quad0<\varepsilon\le cLR}



\textbf{Oracle correspondence.} Applicable interfaces: selected component values and gradients. Algorithms may use only these interfaces and the auxiliary operations explicitly declared public.



\textbf{Guarantee conversion.} A deterministic error bound for every instance also gives the same expected-error bound. Probability guarantees are used only when an explicit conversion to the target contract is recorded.



\textbf{Output rule.} The output must lie in the declared output space and satisfy its feasibility, support, or atomic-hull restrictions. Permission to return an unqueried point does not remove these restrictions.



\textbf{Resource conversion.} Every primitive response is charged. Sample, component, proximal, and communication costs use the units declared for the target setting. Resource coordinates achieved by one algorithm remain a joint vector; coordinatewise minima from different algorithms are not combined.



\textbf{Local argument.} FO\_FINITE\_MAXIMUM, FO\_PROJECTED\_SG, FO\_CERTIFICATE, FO\_NESTED\_FG



\needspace{4\baselineskip}\subsection{B8 - Return the ball center}


\textbf{Admission.} Included in the bracket



\textbf{Direction.} upper



\textbf{View.} Deterministic / Pointwise guarantee



\textbf{Witness or lower-bound scope.} Deterministic witness



\AtlasMath{0}


\textbf{Sources.} \href{https://arxiv.org/pdf/2311.10886}{[3]} \href{https://arxiv.org/pdf/1405.4980}{[4]} \href{https://www2.isye.gatech.edu/~nemirovs/MOR_AccuracyCertificates.pdf}{[5]} \href{https://arxiv.org/pdf/2609.20701}{[6]}



\textbf{Conditions.}
\AtlasMath{\varepsilon\ge LR}




\textbf{Dimension conditions.} \AtlasInlineMath{\varepsilon\ge LR}



\textbf{Oracle correspondence.} Applicable interfaces: selected component values and gradients. Algorithms may use only these interfaces and the auxiliary operations explicitly declared public.



\textbf{Guarantee conversion.} A deterministic error bound for every instance also gives the same expected-error bound. Probability guarantees are used only when an explicit conversion to the target contract is recorded.



\textbf{Output rule.} The output must lie in the declared output space and satisfy its feasibility, support, or atomic-hull restrictions. Permission to return an unqueried point does not remove these restrictions.



\textbf{Resource conversion.} Every primitive response is charged. Sample, component, proximal, and communication costs use the units declared for the target setting. Resource coordinates achieved by one algorithm remain a joint vector; coordinatewise minima from different algorithms are not combined.



\textbf{Local argument.} FO\_FINITE\_MAXIMUM, FO\_PROJECTED\_SG, FO\_CERTIFICATE, FO\_NESTED\_FG



\needspace{4\baselineskip}\subsection{B9 - Accelerated iterative linearization}


\textbf{Admission.} Included in the bracket



\textbf{Direction.} upper



\textbf{View.} Randomized / Expected loss



\textbf{Witness or lower-bound scope.} Deterministic witness



\AtlasMath{Cn(1+\sqrt{B_\varepsilon})}


\textbf{Sources.} \href{https://arxiv.org/pdf/2105.01778}{[1]}



\textbf{Conditions.}
\AtlasMath{0<\varepsilon\le cLR}




\textbf{Dimension conditions.} \AtlasInlineMath{0<\varepsilon\le cLR}



\textbf{Oracle correspondence.} Applicable interfaces: selected component values and gradients. Algorithms may use only these interfaces and the auxiliary operations explicitly declared public.



\textbf{Guarantee conversion.} A deterministic error bound for every instance also gives the same expected-error bound. Probability guarantees are used only when an explicit conversion to the target contract is recorded.



\textbf{Output rule.} The output must lie in the declared output space and satisfy its feasibility, support, or atomic-hull restrictions. Permission to return an unqueried point does not remove these restrictions.



\textbf{Resource conversion.} Every primitive response is charged. Sample, component, proximal, and communication costs use the units declared for the target setting. Resource coordinates achieved by one algorithm remain a joint vector; coordinatewise minima from different algorithms are not combined.



\textbf{Source location.} \href{https://arxiv.org/pdf/2105.01778}{[1]} theorem: Theorem3 Eq.(13), PDF p.12; AppendixA.3 andD.3



\needspace{4\baselineskip}\subsection{B10 - Projected subgradient descent after full value-subgradient simulation using \AtlasInlineMath{n} component calls}


\textbf{Admission.} Included in the bracket



\textbf{Direction.} upper



\textbf{View.} Randomized / Expected loss



\textbf{Witness or lower-bound scope.} Deterministic witness



\AtlasMath{Cn(1+A_\varepsilon^2)}


\textbf{Sources.} \href{https://arxiv.org/pdf/2311.10886}{[3]} \href{https://arxiv.org/pdf/1405.4980}{[4]} \href{https://www2.isye.gatech.edu/~nemirovs/MOR_AccuracyCertificates.pdf}{[5]} \href{https://arxiv.org/pdf/2609.20701}{[6]}



\textbf{Conditions.}
\AtlasMath{0<\varepsilon\le cLR}




\textbf{Dimension conditions.} \AtlasInlineMath{0<\varepsilon\le cLR}



\textbf{Oracle correspondence.} Applicable interfaces: selected component values and gradients. Algorithms may use only these interfaces and the auxiliary operations explicitly declared public.



\textbf{Guarantee conversion.} A deterministic error bound for every instance also gives the same expected-error bound. Probability guarantees are used only when an explicit conversion to the target contract is recorded.



\textbf{Output rule.} The output must lie in the declared output space and satisfy its feasibility, support, or atomic-hull restrictions. Permission to return an unqueried point does not remove these restrictions.



\textbf{Resource conversion.} Every primitive response is charged. Sample, component, proximal, and communication costs use the units declared for the target setting. Resource coordinates achieved by one algorithm remain a joint vector; coordinatewise minima from different algorithms are not combined.



\textbf{Local argument.} FO\_FINITE\_MAXIMUM, FO\_PROJECTED\_SG, FO\_CERTIFICATE, FO\_NESTED\_FG



\needspace{4\baselineskip}\subsection{B11 - Cutting planes after full value-subgradient simulation using \AtlasInlineMath{n} component calls}


\textbf{Admission.} Included in the bracket



\textbf{Direction.} upper



\textbf{View.} Randomized / Expected loss



\textbf{Witness or lower-bound scope.} Deterministic witness



\AtlasMath{Cn[1+d\log(2+A_\varepsilon)]}


\textbf{Sources.} \href{https://www2.isye.gatech.edu/~nemirovs/MOR_AccuracyCertificates.pdf}{[2]}



\textbf{Conditions.}
\AtlasMath{0<\varepsilon\le cLR}




\textbf{Dimension conditions.} \AtlasInlineMath{0<\varepsilon\le cLR}



\textbf{Oracle correspondence.} Applicable interfaces: selected component values and gradients. Algorithms may use only these interfaces and the auxiliary operations explicitly declared public.



\textbf{Guarantee conversion.} A deterministic error bound for every instance also gives the same expected-error bound. Probability guarantees are used only when an explicit conversion to the target contract is recorded.



\textbf{Output rule.} The output must lie in the declared output space and satisfy its feasibility, support, or atomic-hull restrictions. Permission to return an unqueried point does not remove these restrictions.



\textbf{Resource conversion.} Every primitive response is charged. Sample, component, proximal, and communication costs use the units declared for the target setting. Resource coordinates achieved by one algorithm remain a joint vector; coordinatewise minima from different algorithms are not combined.



\textbf{Source location.} \href{https://www2.isye.gatech.edu/~nemirovs/MOR_AccuracyCertificates.pdf}{[2]} theorem: Accuracy certificates; combined with F04 FSG simulation



\needspace{4\baselineskip}\subsection{B12 - Data structure for ball-oracle acceleration}


\textbf{Admission.} Included in the bracket



\textbf{Direction.} upper



\textbf{View.} Randomized / Expected loss



\textbf{Witness or lower-bound scope.} Randomized witness



\AtlasMath{\widetilde{O}(nB_\varepsilon^{1/3}+A_\varepsilon^2)}


\textbf{Sources.} \href{https://arxiv.org/pdf/2311.10886}{[3]} \href{https://arxiv.org/pdf/1405.4980}{[4]} \href{https://www2.isye.gatech.edu/~nemirovs/MOR_AccuracyCertificates.pdf}{[5]} \href{https://arxiv.org/pdf/2609.20701}{[6]}



\textbf{Conditions.}
\AtlasMath{1\le B_\varepsilon\le A_\varepsilon^2,\quad0<\varepsilon\le cLR}




\textbf{Dimension conditions.} \AtlasInlineMath{1\le B_\varepsilon\le A_\varepsilon^2,\quad0<\varepsilon\le cLR}



\textbf{Oracle correspondence.} Applicable interfaces: selected component values and gradients. Algorithms may use only these interfaces and the auxiliary operations explicitly declared public.



\textbf{Guarantee conversion.} After fixing the instance and admissible response law, take expectation over the randomness declared by the target contract and control the underlying error.



\textbf{Output rule.} The output must lie in the declared output space and satisfy its feasibility, support, or atomic-hull restrictions. Permission to return an unqueried point does not remove these restrictions.



\textbf{Resource conversion.} Every primitive response is charged. Sample, component, proximal, and communication costs use the units declared for the target setting. Resource coordinates achieved by one algorithm remain a joint vector; coordinatewise minima from different algorithms are not combined.



\textbf{Local argument.} FO\_FINITE\_MAXIMUM, FO\_PROJECTED\_SG, FO\_CERTIFICATE, FO\_NESTED\_FG



\textbf{Source location.} \href{https://arxiv.org/pdf/2311.10886}{[3]} theorem: Table1; §2.5 iteration/evaluation complexity; Thm8.1; §8.1



\needspace{4\baselineskip}\subsection{B13 - Hard instances for arbitrary randomized component value-gradient algorithms}


\textbf{Admission.} Included in the bracket



\textbf{Direction.} lower



\textbf{View.} Randomized / Expected loss



\textbf{Witness or lower-bound scope.} All randomized adaptive algorithms



\AtlasMath{c[n\min\{A_\varepsilon^{2/3},B_\varepsilon^{1/3}\}+\min\{A_\varepsilon^2,\sqrt{nB_\varepsilon}\}]}


\textbf{Sources.} \href{https://arxiv.org/pdf/2105.01778}{[1]}



\textbf{Conditions.}
\AtlasMath{\varepsilon\le c\min\{LR,\beta R^2\},\quad d\ge C\min\{A_\varepsilon^2,B_\varepsilon\}\log(2+n\min\{A_\varepsilon,B_\varepsilon\})}




\textbf{Dimension conditions.} \AtlasInlineMath{\varepsilon\le c\min\{LR,\beta R^2\},\quad d\ge C\min\{A_\varepsilon^2,B_\varepsilon\}\log(2+n\min\{A_\varepsilon,B_\varepsilon\})}



\textbf{Oracle correspondence.} Applicable interfaces: selected component values and gradients. Lower bounds transfer only through the explicitly recorded simulations from a stronger interface to a weaker one.



\textbf{Guarantee conversion.} The hard-instance quantifiers cover the algorithm class explicitly specified for this bound. Each target setting retains its expected-error or success-probability criterion; a deterministic lower bound is not promoted to a randomized lower bound. Converting the source hard event to an expected-error statement retains the constant tolerance rescaling. Expected running costs additionally require a stopping-time truncation with its constant loss.



\textbf{Output rule.} The output must lie in the declared output space and satisfy its feasibility, support, or atomic-hull restrictions. Permission to return an unqueried point does not remove these restrictions.



\textbf{Resource conversion.} Every primitive response is charged. Sample, component, proximal, and communication costs use the units declared for the target setting. Resource coordinates achieved by one algorithm remain a joint vector; coordinatewise minima from different algorithms are not combined.



\textbf{Source location.} \href{https://arxiv.org/pdf/2105.01778}{[1]} theorem: Theorem3 Eq.(13), PDF p.12; AppendixA.3 andD.3



\needspace{4\baselineskip}\subsection{B14 - Hidden index in arbitrary dimension}


\textbf{Admission.} Included in the bracket



\textbf{Direction.} lower



\textbf{View.} Randomized / Expected loss



\textbf{Witness or lower-bound scope.} All randomized adaptive algorithms



\AtlasMath{cn}


\textbf{Sources.} \href{https://arxiv.org/pdf/2311.10886}{[3]} \href{https://arxiv.org/pdf/1405.4980}{[4]} \href{https://www2.isye.gatech.edu/~nemirovs/MOR_AccuracyCertificates.pdf}{[5]} \href{https://arxiv.org/pdf/2609.20701}{[6]}



\textbf{Conditions.}
\AtlasMath{0<\varepsilon\le cLR}




\textbf{Dimension conditions.} \AtlasInlineMath{0<\varepsilon\le cLR}



\textbf{Oracle correspondence.} Applicable interfaces: selected component values and gradients. Lower bounds transfer only through the explicitly recorded simulations from a stronger interface to a weaker one.



\textbf{Guarantee conversion.} The hard-instance quantifiers cover the algorithm class explicitly specified for this bound. Each target setting retains its expected-error or success-probability criterion; a deterministic lower bound is not promoted to a randomized lower bound.



\textbf{Output rule.} The output must lie in the declared output space and satisfy its feasibility, support, or atomic-hull restrictions. Permission to return an unqueried point does not remove these restrictions.



\textbf{Resource conversion.} Every primitive response is charged. Sample, component, proximal, and communication costs use the units declared for the target setting. Resource coordinates achieved by one algorithm remain a joint vector; coordinatewise minima from different algorithms are not combined.



\textbf{Local argument.} FO\_FINITE\_MAXIMUM, FO\_PROJECTED\_SG, FO\_CERTIFICATE, FO\_NESTED\_FG



\needspace{4\baselineskip}\subsection{B15 - Replication of a one-dimensional smooth hard objective}


\textbf{Admission.} Included in the bracket



\textbf{Direction.} lower



\textbf{View.} Randomized / Expected loss



\textbf{Witness or lower-bound scope.} All randomized adaptive algorithms



\AtlasMath{c\log(2+\min\{LR,\beta R^2\}/\varepsilon)}


\textbf{Sources.} \href{https://arxiv.org/pdf/2311.10886}{[3]} \href{https://arxiv.org/pdf/1405.4980}{[4]} \href{https://www2.isye.gatech.edu/~nemirovs/MOR_AccuracyCertificates.pdf}{[5]} \href{https://arxiv.org/pdf/2609.20701}{[6]}



\textbf{Conditions.}
\AtlasMath{0<\varepsilon\le c\min\{LR,\beta R^2\}}




\textbf{Dimension conditions.} \AtlasInlineMath{0<\varepsilon\le c\min\{LR,\beta R^2\}}



\textbf{Oracle correspondence.} Applicable interfaces: selected component values and gradients. Lower bounds transfer only through the explicitly recorded simulations from a stronger interface to a weaker one.



\textbf{Guarantee conversion.} The hard-instance quantifiers cover the algorithm class explicitly specified for this bound. Each target setting retains its expected-error or success-probability criterion; a deterministic lower bound is not promoted to a randomized lower bound.



\textbf{Output rule.} The output must lie in the declared output space and satisfy its feasibility, support, or atomic-hull restrictions. Permission to return an unqueried point does not remove these restrictions.



\textbf{Resource conversion.} Every primitive response is charged. Sample, component, proximal, and communication costs use the units declared for the target setting. Resource coordinates achieved by one algorithm remain a joint vector; coordinatewise minima from different algorithms are not combined.



\textbf{Local argument.} FO\_FINITE\_MAXIMUM, FO\_PROJECTED\_SG, FO\_CERTIFICATE, FO\_NESTED\_FG



\needspace{4\baselineskip}\subsection{B16 - Recover all affine component formulas}


\textbf{Admission.} Included in the bracket



\textbf{Direction.} upper



\textbf{View.} Randomized / Expected loss



\textbf{Witness or lower-bound scope.} Deterministic witness



\AtlasMath{n}


\textbf{Sources.} \href{https://arxiv.org/pdf/2311.10886}{[3]} \href{https://arxiv.org/pdf/1405.4980}{[4]} \href{https://www2.isye.gatech.edu/~nemirovs/MOR_AccuracyCertificates.pdf}{[5]} \href{https://arxiv.org/pdf/2609.20701}{[6]}



\textbf{Conditions.}
\AtlasMath{\beta=0,\quad0<\varepsilon\le cLR}




\textbf{Dimension conditions.} \AtlasInlineMath{\beta=0,\quad0<\varepsilon\le cLR}



\textbf{Oracle correspondence.} Applicable interfaces: selected component values and gradients. Algorithms may use only these interfaces and the auxiliary operations explicitly declared public.



\textbf{Guarantee conversion.} A deterministic error bound for every instance also gives the same expected-error bound. Probability guarantees are used only when an explicit conversion to the target contract is recorded.



\textbf{Output rule.} The output must lie in the declared output space and satisfy its feasibility, support, or atomic-hull restrictions. Permission to return an unqueried point does not remove these restrictions.



\textbf{Resource conversion.} Every primitive response is charged. Sample, component, proximal, and communication costs use the units declared for the target setting. Resource coordinates achieved by one algorithm remain a joint vector; coordinatewise minima from different algorithms are not combined.



\textbf{Local argument.} FO\_FINITE\_MAXIMUM, FO\_PROJECTED\_SG, FO\_CERTIFICATE, FO\_NESTED\_FG



\needspace{4\baselineskip}\subsection{B17 - Return the ball center}


\textbf{Admission.} Included in the bracket



\textbf{Direction.} upper



\textbf{View.} Randomized / Expected loss



\textbf{Witness or lower-bound scope.} Deterministic witness



\AtlasMath{0}


\textbf{Sources.} \href{https://arxiv.org/pdf/2311.10886}{[3]} \href{https://arxiv.org/pdf/1405.4980}{[4]} \href{https://www2.isye.gatech.edu/~nemirovs/MOR_AccuracyCertificates.pdf}{[5]} \href{https://arxiv.org/pdf/2609.20701}{[6]}



\textbf{Conditions.}
\AtlasMath{\varepsilon\ge LR}




\textbf{Dimension conditions.} \AtlasInlineMath{\varepsilon\ge LR}



\textbf{Oracle correspondence.} Applicable interfaces: selected component values and gradients. Algorithms may use only these interfaces and the auxiliary operations explicitly declared public.



\textbf{Guarantee conversion.} A deterministic error bound for every instance also gives the same expected-error bound. Probability guarantees are used only when an explicit conversion to the target contract is recorded.



\textbf{Output rule.} The output must lie in the declared output space and satisfy its feasibility, support, or atomic-hull restrictions. Permission to return an unqueried point does not remove these restrictions.



\textbf{Resource conversion.} Every primitive response is charged. Sample, component, proximal, and communication costs use the units declared for the target setting. Resource coordinates achieved by one algorithm remain a joint vector; coordinatewise minima from different algorithms are not combined.



\textbf{Local argument.} FO\_FINITE\_MAXIMUM, FO\_PROJECTED\_SG, FO\_CERTIFICATE, FO\_NESTED\_FG



\needspace{4\baselineskip}\subsection{B18 - Randomized affine-minorant local solver and constrained acceleration}


\textbf{Admission.} Included in the bracket



\textbf{Direction.} upper



\textbf{View.} Randomized / 2/3 success / normalized global regularity



\textbf{Witness or lower-bound scope.} Randomized witness



\AtlasMath{C(n\varepsilon^{-1/3}+\sqrt{n/\varepsilon})\log^6(en/\varepsilon)}


\textbf{Sources.} \href{https://github.com/WU-Wendao/Colombo/blob/main/colombo_sensorium/Complexity_altas/data/sources/colombo/A2/manuscript.pdf}{[7]}



\textbf{Conditions.} \AtlasInlineMath{n\ge1} and \AtlasInlineMath{0<\varepsilon<1}; the complexity takes the worst case over all finite dimensions.; The constructed algorithm actually succeeds with probability at least \AtlasInlineMath{0.99}. Every full pass costs \AtlasInlineMath{n} queries.; The normalized global-regularity view is distinct from assuming regularity only near the ball. The comparison concerns oracle calls, not arithmetic runtime.



\textbf{Dimension conditions.} The upper bound holds for every finite d. The lower bound is existential in dimension, which may depend on n and epsilon; A2 Proposition prop:lower and Section sec:lower do not give an explicit dimensional formula.



\textbf{Constants.} Universal positive constants.



\textbf{Logarithmic factors.} Only the displayed sixth power of \AtlasInlineMath{\log(en/\varepsilon)}.



\textbf{Source location.} \href{https://github.com/WU-Wendao/Colombo/blob/main/colombo_sensorium/Complexity_altas/data/sources/colombo/A2/manuscript.pdf}{[7]} Theorem thm:main, eq:upper; Sections sec:local to sec:combine



Randomized affine-minorant local solver and constrained acceleration. Unrestricted adaptive randomized algorithms



The constructed algorithm actually succeeds with probability at least \AtlasInlineMath{0.99}. Every full pass costs \AtlasInlineMath{n} queries.



\needspace{4\baselineskip}\subsection{B19 - Carmon--Jambulapati--Jin--Sidford randomized local-oracle lower bound, with the stated normalization and terminal-query reduction}


\textbf{Admission.} Included in the bracket



\textbf{Direction.} lower



\textbf{View.} Randomized / 2/3 success / normalized global regularity



\textbf{Witness or lower-bound scope.} randomized\_hard\_family



\AtlasMath{c(n\varepsilon^{-1/3}+\sqrt{n/\varepsilon})}


\textbf{Sources.} \href{https://github.com/WU-Wendao/Colombo/blob/main/colombo_sensorium/Complexity_altas/data/sources/colombo/A2/manuscript.pdf}{[7]}



\textbf{Conditions.} \AtlasInlineMath{n\ge1} and \AtlasInlineMath{0<\varepsilon<1}; the complexity takes the worst case over all finite dimensions.; The lower bound is not asserted for each fixed dimension.; The normalized global-regularity view is distinct from assuming regularity only near the ball. The comparison concerns oracle calls, not arithmetic runtime.



\textbf{Dimension conditions.} The upper bound holds for every finite d. The lower bound is existential in dimension, which may depend on n and epsilon; A2 Proposition prop:lower and Section sec:lower do not give an explicit dimensional formula.



\textbf{Constants.} Universal positive constants.



\textbf{Logarithmic factors.} Only the displayed sixth power of \AtlasInlineMath{\log(en/\varepsilon)}.



\textbf{Source location.} \href{https://github.com/WU-Wendao/Colombo/blob/main/colombo_sensorium/Complexity_altas/data/sources/colombo/A2/manuscript.pdf}{[7]} Theorem thm:main; Proposition prop:lower; Lemma lem:constant-lower



Carmon--Jambulapati--Jin--Sidford randomized local-oracle lower bound, with the stated normalization and terminal-query reduction. All unrestricted adaptive randomized algorithms



The lower bound is not asserted for each fixed dimension.



\needspace{7\baselineskip}\section{Maintenance appendix: Native proof extensions}


Read-only presentation of the stored native record. Integrity hashes identify source values; this appendix grants no new theorem, source-read credit, review or approval.



Native pointer: /proofs; canonical SHA-256: d31d35a2d645dbd785b9f519\allowbreak{}a36a1d8e323166a1b8accabc\allowbreak{}d9463a380e5bcfab



\needspace{4\baselineskip}\subsection{Native proof: FO\_FINITE\_MAXIMUM}


Native pointer: /proofs/0; canonical SHA-256: 821d35bc5b288862cb88bf4a\allowbreak{}d41a5f3109c1257a28db9ceb\allowbreak{}3fdffe0b29fbb4af



\needspace{4\baselineskip}\subsubsection{Native field: title}


Native pointer: /proofs/0/title; canonical SHA-256: db63ec3d90ae478a08f56909\allowbreak{}7f01f711e22ce793ad293a31\allowbreak{}a33aefb4bb377931


Finite maxima, hidden offsets, and budget conversion


\needspace{4\baselineskip}\subsubsection{Native field: body}


Native pointer: /proofs/0/body; canonical SHA-256: ed3e5d317e854275eba7a207\allowbreak{}1ffd868610b850c9e55c0a4a\allowbreak{}6366fe906efa6af2


With component value\char45{}gradient access, query all \$n\$ components, select an active one, and use its gradient as a subgradient of \$F=\textbackslash{}max\_i f\_i\$. This transfers subgradient and certificate algorithms at cost \$n\$ per full response. The maximum is \$L\$\char45{}Lipschitz; the center has gap at most \$LR\$.

For an \$\textbackslash{}Omega(n)\$ lower bound, hide an index \$J\$. Its component is \$LR+Lx\_1\$ or \$LR\char45{}Lx\_1\$; every other component is zero. The hidden component is nonnegative on the ball and dominates the maximum. Before it is queried, the two signs have average loss \$LR\$ for any common output. A fixed budget hits \$J\$ with probability at most \$T/n\$. Thus \$\textbackslash{}Omega(n)\$ calls are necessary at nontrivial accuracy, for arbitrary randomized algorithms.

When \$\textbackslash{}beta=0\$, each component is affine. One value\char45{}gradient query reconstructs its formula, so \$n\$ calls and free internal optimization solve the problem exactly. Replicating a smooth one\char45{}dimensional hard function in all components gives the logarithmic lower term with scale \$\textbackslash{}min\textbackslash{}\{LR,\textbackslash{}beta R\textasciicircum{}2\textbackslash{}\}\$. Combine it with the hidden\char45{}index term by taking a maximum, not a product.

With component gradients alone, relative additive offsets are invisible. The pair \$Lx\$, \$\char45{}Lx+c\$ has the same gradients for all \$c\$, but its maximum has different minimizers. Two offsets with separated interior minimizers give a positive worst\char45{}case error for every common output distribution, independently of the query budget.

For an imported upper bound with a constant\char45{}probability joint event of acceptable error and budget, truncate each run, repeat independently, and evaluate all components at every candidate to choose the best exact maximum. Feasible outputs have gap at most \$2LR\$. Reducing failure to \$\textbackslash{}epsilon/(4LR)\$ therefore gives expected error at most \$\textbackslash{}epsilon\$. Repetition and validation calls are charged, producing the recorded logarithmic overhead.


\needspace{4\baselineskip}\subsubsection{Native field: scope}


Native pointer: /proofs/0/scope; canonical SHA-256: dd9d9d36aa8db4b6692984bc\allowbreak{}92345d4dfc1747f2c6ba246c\allowbreak{}5a5e376b6554a7a3


Component values and gradients for simulation; component gradients alone for the offset impossibility.


\needspace{4\baselineskip}\subsubsection{Native field: limitations}


Native pointer: /proofs/0/limitations; canonical SHA-256: 30096e4bead8971af01075b1\allowbreak{}2f565e8fac1ec0e7e03ee997\allowbreak{}0ce1beab0154e465


Additive offsets and the guarantee/budget conversion must be kept explicit.


\needspace{4\baselineskip}\subsubsection{Native field: algorithm\_scope}


Native pointer: /proofs/0/algorithm\_scope; canonical SHA-256: 544a2cce941e4af69392892b\allowbreak{}7bc87f6122f46b22f4336855\allowbreak{}1d9943e2454a1c33


Deterministic simulation upper witnesses; randomized hidden\char45{}index lower bounds; randomized budget conversion when needed.


\needspace{4\baselineskip}\subsubsection{Native field: source\_refs}


Native pointer: /proofs/0/source\_refs; canonical SHA-256: 0929d2aa40221e9b3ae9a1bd\allowbreak{}6a88a31b4f860912081307b6\allowbreak{}a1bd8c50cae83916



\needspace{4\baselineskip}\subsubsection{Native item 1}


Native pointer: /proofs/0/source\_refs/0; canonical SHA-256: 4d259391de34ed28606e7ded\allowbreak{}aa4bd7420cdca35d4374466e\allowbreak{}c09a76171549cfc9


SO\_SRC\_1ccf639c7d0ff459


\needspace{4\baselineskip}\subsection{Native proof: FO\_PROJECTED\_SG}


Native pointer: /proofs/1; canonical SHA-256: 2097db46c9b028dabab6fcc3\allowbreak{}e9b82ce6016c473a751004b4\allowbreak{}2f2697573084d90c



\needspace{4\baselineskip}\subsubsection{Native field: title}


Native pointer: /proofs/1/title; canonical SHA-256: 96370c58f6530dbf1c90377d\allowbreak{}7fc9b80f22df86a70b71b137\allowbreak{}163b08ba7b7d4d10


Projected and weighted subgradient bounds


\needspace{4\baselineskip}\subsubsection{Native field: body}


Native pointer: /proofs/1/body; canonical SHA-256: df6b6206889b23cb4b2b33b6\allowbreak{}4b228616a05fe162f308b80e\allowbreak{}a085037fc7cda3e6


Projection onto a public closed convex set gives
\$\$
\textbackslash{}|x\_\{t+1\}\char45{}x\textasciicircum{}\textbackslash{}star\textbackslash{}|\textasciicircum{}2\textbackslash{}le\allowbreak{}\textbackslash{}|x\_t\char45{}x\textasciicircum{}\textbackslash{}star\textbackslash{}|\textasciicircum{}2
\char45{}2\textbackslash{}eta\_t\textbackslash{}langle g\_t,x\_t\char45{}x\textasciicircum{}\textbackslash{}star\textbackslash{}rangle+\textbackslash{}\allowbreak{}eta\_t\textasciicircum{}2\textbackslash{}|g\_t\textbackslash{}|\textasciicircum{}2.
\$\$
For convex \$f\$, initial distance at most \$R\$, and conditional second moment at most \$V\$, the constant step \$R/\textbackslash{}sqrt\{VK\}\$ and averaged output have expected gap at most \$R\textbackslash{}sqrt\{V/K\}\$. Exact subgradients correspond to \$V=L\textasciicircum{}2\$.

Under \$\textbackslash{}mu\$\char45{}strong convexity, use \$\textbackslash{}eta\_t=2/[\textbackslash{}mu(t+1)]\$. With \$a\_t=\textbackslash{}mathbb E\textbackslash{}|x\_t\char45{}x\textasciicircum{}\textbackslash{}star\textbackslash{}|\textasciicircum{}2\$,
\$\$
\textbackslash{}mathbb E[f(x\_t)\char45{}f\textasciicircum{}\textbackslash{}star]\textbackslash{}le\{\textbackslash{}mu(t\char45{}1)\textbackslash{}over4\}a\_t
\char45{}\{\textbackslash{}mu(t+1)\textbackslash{}over4\}a\_\{t+1\}\allowbreak{}+\{V\textbackslash{}over\textbackslash{}mu(t+1)\}.
\$\$
Multiply by \$t\$, sum, and use convexity. The weighted output
\$\textbackslash{}sum\_\{t=1\}\textasciicircum{}K2t\textbackslash{},x\_t/[K(K+1)]\$ has expected gap at most \$2V/[\textbackslash{}mu(K+1)]\$. Neither output rule requires private randomness. Apply norm conversion before these inequalities, preserving all dimension factors.


\needspace{4\baselineskip}\subsubsection{Native field: scope}


Native pointer: /proofs/1/scope; canonical SHA-256: aefb58efc6774fa2bbc4a16c\allowbreak{}0a22ce278c301d969fa7e33f\allowbreak{}7eb67f48d0b390dd


Public Euclidean projection; unbiased conditional subgradient mean and a total second\char45{}moment bound, or exact subgradients.


\needspace{4\baselineskip}\subsubsection{Native field: limitations}


Native pointer: /proofs/1/limitations; canonical SHA-256: efe5c4306e0e4faf5d2c8e1f\allowbreak{}00f38151f773bb9855ee8d9a\allowbreak{}ffde1fa5b9c9771f


Norm\char45{}conversion factors must be applied explicitly. No noisy cutting\char45{}plane result follows from this recursion.


\needspace{4\baselineskip}\subsubsection{Native field: algorithm\_scope}


Native pointer: /proofs/1/algorithm\_scope; canonical SHA-256: e1a31a03176516390daf9011\allowbreak{}669c9af69332f5c849e08101\allowbreak{}3b0de5e2158ecd13


Deterministic policies and deterministic averages; randomness may come from observations.


\needspace{4\baselineskip}\subsection{Native proof: FO\_CERTIFICATE}


Native pointer: /proofs/2; canonical SHA-256: 21979664af3d9c988ec0425d\allowbreak{}315ce0b043b36ff9d545026e\allowbreak{}d3322eb72c59ff61



\needspace{4\baselineskip}\subsubsection{Native field: title}


Native pointer: /proofs/2/title; canonical SHA-256: a3ebc976e1b51b6b06145bd0\allowbreak{}666cc8cefdf1c282c07da232\allowbreak{}38d9b82f3a05b3cb


Subgradient certificates and centroid cuts


\needspace{4\baselineskip}\subsubsection{Native field: body}


Native pointer: /proofs/2/body; canonical SHA-256: 72d8341a17456c969432c1e8\allowbreak{}06fc1bc7998d07720b924200\allowbreak{}71af653c600ce2eb


Let \$Q\textbackslash{}subset\textbackslash{}mathbb R\textasciicircum{}d\$ be public, compact, convex, and full dimensional, with diameter at most \$D\$ in the stated norm. Suppose all returned subgradients have dual norm at most \$L\$. Query \$g\_i\textbackslash{}in\textbackslash{}partial f(x\_i)\$ at the centroid of the current body and retain \$\textbackslash{}langle g\_i,x\_i\char45{}x\textbackslash{}rangle\textbackslash{}ge0\$. The centroid halfspace theorem bounds the retained volume by \$(1\char45{}e\textasciicircum{}\{\char45{}1\})\textasciicircum{}T\textbackslash{}operatorname\{vol\}(Q)\$.

The computable certificate is
\$\$
\textbackslash{}Gamma\_T=\textbackslash{}min\_\{\textbackslash{}lambda\textbackslash{}i\allowbreak{}n\textbackslash{}Delta\_T\}\textbackslash{}max\_\{x\textbackslash{}in Q\}\textbackslash{}sum\_i\textbackslash{}lambda\_i\textbackslash{}langle g\_i,x\_i\char45{}x\textbackslash{}rangle
=\textbackslash{}max\_\{x\textbackslash{}in Q\}\textbackslash{}min\_i\textbackslash{}langle g\_i,x\_i\char45{}x\textbackslash{}rangle.
\$\$
The equality is convex minimax. Subgradient inequalities and convexity imply
\$f(\textbackslash{}sum\_i\textbackslash{}lambda\_ix\_i)\char45{}\textbackslash{}\allowbreak{}min\_Qf\textbackslash{}le\textbackslash{}Gamma\_T\$.
If \$\textbackslash{}Gamma\_T>\textbackslash{}epsilon\$, choose \$z\textbackslash{}in Q\$ with all cut margins greater than \$\textbackslash{}epsilon\$. The copy \$(1\char45{}a)z+aQ\$, where \$a=\textbackslash{}epsilon/[2(\textbackslash{}epsilon+LD)]\$, survives every cut. Its volume is \$a\textasciicircum{}d\textbackslash{}operatorname\{vol\}(Q)\$, contradicting centroid shrinkage once \$T>C d\textbackslash{}log(2+LD/\textbackslash{}epsilon)\$. If the body loses volume earlier, the same positive\char45{}volume argument already certifies the error. No function values or inscribed\char45{}ball estimate are required.

For \$Q=RB\_p\textasciicircum{}d\$ and \$\textbackslash{}ell\_q\$\char45{}Lipschitz objectives, use \$D=2Rd\textasciicircum{}\{(1/q\char45{}1/p)\_+\}\$. Returning the center has error at most \$LRd\textasciicircum{}\{(1/q\char45{}1/p)\_+\}\$, giving the coarse zero\char45{}query regime. Exact centroid, minimax, and public geometric calculations are uncharged internal operations.

For a deterministic selector with total movement at most \$W(T)\$ on nested convex sets, instead impose deep cuts of margin \$2LW(T)/T\$. If all \$T\$ sets are nonempty, every move has length at least \$2W(T)/T\$, contradicting the movement bound. At the first empty intersection, minimax gives error at most \$2LW(T)/T\$. Substitution of the cited stable\char45{}movement theorem yields the recorded nondual upper bound. Its deterministic information scope is distinct from the scope of randomized implementations.


\needspace{4\baselineskip}\subsubsection{Native field: scope}


Native pointer: /proofs/2/scope; canonical SHA-256: 92d674be8fd361e7686d71fa\allowbreak{}4f979260ea491c9b01566433\allowbreak{}4a0ec353e8059b63


One full subgradient is sufficient. Public compact convex geometry and exact internal geometric computation are available.


\needspace{4\baselineskip}\subsubsection{Native field: limitations}


Native pointer: /proofs/2/limitations; canonical SHA-256: 3c41670d63fc68cbd9a2b20c\allowbreak{}5bdffe5abdf747751c9b855a\allowbreak{}1dcef06421a036cd


The stable\char45{}selector branch imports its movement theorem and does not claim efficient deterministic arithmetic time.


\needspace{4\baselineskip}\subsubsection{Native field: algorithm\_scope}


Native pointer: /proofs/2/algorithm\_scope; canonical SHA-256: 42d065ebafc108a64e9cee96\allowbreak{}4510776065fd1792b4344749\allowbreak{}b3f72a002ccd46aa


Deterministic upper\char45{}bound witness; valid as an upper bound for randomized views.


\needspace{4\baselineskip}\subsubsection{Native field: source\_refs}


Native pointer: /proofs/2/source\_refs; canonical SHA-256: 6920d63b36e4e336e449c0c5\allowbreak{}07724c9a8e022e1682393587\allowbreak{}8f5a8459ea8654b1



\needspace{4\baselineskip}\subsubsection{Native item 1}


Native pointer: /proofs/2/source\_refs/0; canonical SHA-256: f636ed5cb2c3078753fe74a8\allowbreak{}d5e61f525482a9f5b53b6ac3\allowbreak{}4ca8401dd49dc086


SO\_SRC\_a6e5b4d495f30052


\needspace{4\baselineskip}\subsubsection{Native item 2}


Native pointer: /proofs/2/source\_refs/1; canonical SHA-256: 9daedb2ac95e8088053c7f30\allowbreak{}d00e8b484004c953b15791fc\allowbreak{}665f3978350e6bba


SO\_SRC\_c314bae252a3b929


\needspace{4\baselineskip}\subsubsection{Native item 3}


Native pointer: /proofs/2/source\_refs/2; canonical SHA-256: f49cfd2461dd4b0132f4f0fa\allowbreak{}2a90e3f29883781b5ad0bff1\allowbreak{}9390f24edc0bff0c


SO\_SRC\_88e665bbd463b203


\needspace{4\baselineskip}\subsection{Native proof: FO\_NESTED\_FG}


Native pointer: /proofs/3; canonical SHA-256: b78b6a3171d3b379278a8ff8\allowbreak{}bd423f5ef452d85e374e045f\allowbreak{}39f303b55c8d770b



\needspace{4\baselineskip}\subsubsection{Native field: title}


Native pointer: /proofs/3/title; canonical SHA-256: 3ef610b063d3215728ad6bd9\allowbreak{}ca0fb250506c62bb1e79d0a1\allowbreak{}f1712a9223f10515


A one\char45{}dimensional lower bound with exact values and gradients


\needspace{4\baselineskip}\subsubsection{Native field: body}


Native pointer: /proofs/3/body; canonical SHA-256: 8c01c6437e24e5476cf2586f\allowbreak{}5db0d4209f21418f9c6af775\allowbreak{}1dd6cd66e78e7c13


Fix \$1<\textbackslash{}kappa=\textbackslash{}beta/\textbackslash{}mu\$. Choose a fixed smooth \$\textbackslash{}psi\$, equal to \$\char45{}x\$ on \$[\char45{}1/2,1/2]\$ and zero outside \$[\char45{}1,1]\$, and let \$C\_\textbackslash{}psi=\textbackslash{}max\textbackslash{}\{1,\textbackslash{}|\textbackslash{}psi\char39{}\char39{}\textbackslash{}|\_\textbackslash{}infty\textbackslash{}\}\$. Set
\$\$
\textbackslash{}beta\_0=(\textbackslash{}beta+\textbackslash{}mu)/2,\textbackslash{}quad a=\{\textbackslash{}kappa\char45{}1\textbackslash{}over8C\_\textbackslash{}psi(\allowbreak{}\textbackslash{}kappa+1)\},\textbackslash{}quad
\textbackslash{}eta=a/4,\textbackslash{}quad \textbackslash{}rho\_0=\{(1\char45{}\textbackslash{}eta)\textbackslash{}over a\}\textbackslash{}sqrt\{\textbackslash{}Delta/\textbackslash{}beta\}.
\$\$
For independent uniform signs \$b\_i\$, define \$u\_0=0\$,
\$u\_i=u\_\{i\char45{}1\}+b\_i a\textbackslash{}rho\_\{i\char45{}1\}\$, \$\textbackslash{}rho\_i=\textbackslash{}eta\textbackslash{}rho\_\{i\char45{}1\}\$, and
\$\$
f\_b(x)=\textbackslash{}beta\_0x\textasciicircum{}2/2+
\textbackslash{}beta\_0\textbackslash{}sum\_\{i=1\}\textasciicircum{}N b\_i a\textbackslash{}rho\_\{i\char45{}1\}\textasciicircum{}2
\textbackslash{}psi((x\char45{}u\_\{i\char45{}1\})/\textbackslash{}rho\_\{i\char45{}1\}).
\$\$
The two child intervals \$I\_i=[u\_i\char45{}\textbackslash{}rho\_i,u\_i+\textbackslash{}rho\_i]\$ are disjoint and lie inside the central half of the parent. Second derivatives of successive bumps have disjoint shell supports. Thus \$\textbackslash{}mu\textbackslash{}le f\_b\char39{}\char39{}\textbackslash{}le\textbackslash{}beta\$, and \$f\_b\char39{}\char39{}\textbackslash{}ge7\textbackslash{}beta/16\$. Within \$I\_i\$, the partial sum through level \$i\$ is a quadratic centered at \$u\_i\$. The final minimizer is \$u\_N\$, with \$|u\_N|\textbackslash{}le\textbackslash{}sqrt\{\textbackslash{}Delta/\textbackslash{}beta\}\$ and initial gap at most \$\textbackslash{}Delta/2\$.

A stronger oracle reveals the next bit when the query lies in the current unresolved interval, continuing only while it lies in the selected child. Unrevealed bumps vanish near the query, so the revealed prefix computes the exact value and derivative. Conditional on history, a fixed point belongs to at most one child; one query reveals at most \$\textbackslash{}sum\_\{j\textbackslash{}ge1\}2\textasciicircum{}\{\char45{}(j\char45{}1)\}=2\$ new bits in expectation.

Treat the final output as one more such query. If \$|\textbackslash{}widehat x\char45{}u\_N|<\textbackslash{}rho\_N\$, all \$N\$ bits are revealed. Its probability is at most \$2(T+1)/N\$. Otherwise the loss is at least
\$c\textbackslash{}beta\textbackslash{}rho\_N\textasciicircum{}2=c\char39{}\textbackslash{}Delta\textbackslash{}eta\textasciicircum{}\{2(N\char45{}1)\}\$.
Choose \$N\$ at the accuracy scale and use opposite translations to exclude zero calls. For expected error and constant\char45{}probability accuracy,
\$\$
T\textbackslash{}ge c\textbackslash{}left[1+\{\textbackslash{}log(\textbackslash{}Delta/\textbackslash{}epsilon)\textbackslash{}over
\textbackslash{}log(C\textbackslash{}kappa/(\textbackslash{}kappa\char45{}1))\}\textbackslash{}right].
\$\$
Queries may be arbitrary, randomized, adaptive, and unbounded.

Fixing \$\textbackslash{}kappa=2\$, \$\textbackslash{}Delta=c\textbackslash{}beta R\textasciicircum{}2\$, gives the initial\char45{}distance logarithmic bound. The gradient version uses
\$|f\_b\char39{}(x)|\textbackslash{}ge(7\textbackslash{}beta/16)|x\char45{}u\_N|\$.
Add a public quadratic to embed in higher dimension. For composite minimization, apply the construction on total curvature interval \$[\textbackslash{}mu,\textbackslash{}mu+\textbackslash{}beta]\$, with public \$h(x)=\textbackslash{}mu x\textasciicircum{}2/2\$; the unknown remainder is convex and \$\textbackslash{}beta\$\char45{}smooth, and every proximal response for \$h\$ is independent of the hidden bits.


\needspace{4\baselineskip}\subsubsection{Native field: scope}


Native pointer: /proofs/3/scope; canonical SHA-256: 9775156c58630feb2557b65d\allowbreak{}993be46ef31e5b08c31436e8\allowbreak{}7aeedcf91b079fff


Exact values and gradients; unbounded real\char45{}line queries, with higher\char45{}dimensional embeddings.


\needspace{4\baselineskip}\subsubsection{Native field: limitations}


Native pointer: /proofs/3/limitations; canonical SHA-256: a4bee836784937754fb75497\allowbreak{}2e6a463a835e8d6b43165c89\allowbreak{}06fb2a1f5aefcf46


Curvature, initialization, and nontrivial\char45{}accuracy restrictions remain explicit; the composite transfer lower\char45{}bounds unknown\char45{}gradient calls.


\needspace{4\baselineskip}\subsubsection{Native field: algorithm\_scope}


Native pointer: /proofs/3/algorithm\_scope; canonical SHA-256: ec0da73a9cb6b357abfe64f2\allowbreak{}f10f9c420bd6fc8a395b3bce\allowbreak{}0dd79883015a60f5


Lower bounds against arbitrary randomized adaptive algorithms, including unqueried output.


\needspace{7\baselineskip}\section{References}


[1] Yair Carmon, Arun Jambulapati, Yujia Jin and Aaron Sidford. Thinking Inside the Ball: Near-Optimal Minimization of the Maximal Loss. COLT 2021, full arXiv version; Theorem 2 and Eq.(9), p.9; Theorem 3 and Eq.(13), p.12; Appendices A.3 and D.3.



\url{https://arxiv.org/pdf/2105.01778}



[2] Arkadi Nemirovski; Shmuel Onn; Uriel G. Rothblum. Accuracy Certificates for Computational Problems with Convex Structure. 2010; Proposition 2.1(ii), PDF p.5; §4.1.2, PDF p.11.MOR 35(1):52-78..



\url{https://www2.isye.gatech.edu/~nemirovs/MOR_AccuracyCertificates.pdf}



[3] Yair Carmon, Arun Jambulapati, Yujia Jin and Aaron Sidford. A Whole New Ball Game: A Primal Accelerated Method for Matrix Games and Minimizing the Maximum of Smooth Functions. 2023 preprint; STOC 2024; Table 1; Section 2.5; Definitions 8.1-8.2; Theorem 8.1 and Section 8.1.



\url{https://arxiv.org/pdf/2311.10886}



[4] Sébastien Bubeck. Convex Optimization: Algorithms and Complexity. 2015.



\url{https://arxiv.org/pdf/1405.4980}



[5] Arkadi Nemirovski; Shmuel Onn; Uriel G. Rothblum. Accuracy Certificates for Computational Problems with Convex Structure. 2010.



\url{https://www2.isye.gatech.edu/~nemirovs/MOR_AccuracyCertificates.pdf}



[6] David Martínez-Rubio; Cristóbal Guzmán. Stable Movement for Nondual Lipschitz Convex Optimization: Efficiency and Nearly Optimal Oracle Rates. 2026-09-17; preprintv1.



\url{https://arxiv.org/pdf/2609.20701}



[7] Colombo\_A2. Near-Optimal Component-Oracle Complexity for Smooth Maximum Minimization. 2026. A2\_randomized\_affine\_min\allowbreak{}orant\_matching.tex.



\url{https://github.com/WU-Wendao/Colombo/blob/main/colombo_sensorium/Complexity_altas/data/sources/colombo/A2/manuscript.pdf}

\end{document}
